%=====================================================================
\chapter{The Toolkit: Setting It Up and Keeping It}\label{app:toolkit}
%=====================================================================
\normalfigures

This appendix is the whole of the course's setup. It is short, and that is the
point: the longest-running complaint about database courses is that the first
fortnight goes on installing something. Here there is nothing to install.

\section{What You Need}

\rowcolors{2}{white}{lightbg}
\begin{longtable}{@{}L{3.4cm}L{4.2cm}L{6.8cm}@{}}
\toprule
\rowcolor{primary!15}
\textbf{Thing} & \textbf{Version used here} & \textbf{How to get it} \\
\midrule
\endhead
\textbf{SQLite} & 3.51.2 & \textbf{You already have it.} It ships inside
Python, inside most operating systems, and inside a great deal of software you
already run \\
A terminal & any & Already present \\
Python & 3.13 & Only for the measurement harnesses in \texttt{code/}. No
chapter's \emph{content} needs it \\
A text editor & any & For writing \texttt{.sql} files \\
\bottomrule
\end{longtable}

\begin{conceptbox}[Check, in one line, before you read any further]
\begin{lstlisting}[language=bash]
python -c "import sqlite3; print(sqlite3.sqlite_version)"
\end{lstlisting}

If that prints a version number --- any version from 3.37 onwards --- you have
everything this course needs. If you also have the \texttt{sqlite3} command on
your path, so much the better; if not, every lab works through Python's
built-in shell instead, and \chref{firstdatabase} shows both.
\end{conceptbox}

\section{Why SQLite, and What It Costs You}

\begin{alertbox}[This is a curriculum decision, not a shortcut]
\textbf{HEC names no database product in any edition of this curriculum.} The
2017 and 2023 outlines say a student should use SQL ``in any DBMS''; the 2025
outcomes say ``popular database management systems''. There is no edition in
which choosing SQLite departs from the published syllabus.

\smallskip
What it buys is the first fortnight. There is no server to start, no port, no
password, no user account, no configuration file and nothing to go wrong
before the subject begins. A database is one file. You can copy it, email it,
put it in version control, delete it and start again --- and in \chref{normalforms}
you will do exactly that.

\smallskip
\textbf{What it costs is stated here and again in \chref{project}}, because a
tool's limits are part of what you are learning:

\begin{tight}
  \item \textbf{No users, no \sqlkw{GRANT}, no \sqlkw{REVOKE}.} SQLite has no
        access control of any kind: whoever can open the file can do anything
        to it. \chref{project} writes the SQL a server DBMS would use and runs
        it nowhere.
  \item \textbf{No stored procedures.} Triggers (\chref{views}) carry the
        argument about putting logic in the database.
  \item \textbf{No row-level locking and no isolation levels.} SQLite lets one
        writer at a time and makes the others wait. \chref{transactions} can
        therefore show you a lost update prevented, a busy timeout and a failed
        upgrade from a deferred transaction --- but a genuine two-row deadlock
        is drawn rather than run.
\end{tight}

\smallskip
Every one of those three is a thing a \emph{server} adds, and knowing which
three they are is more useful than having had them switched on all semester
without noticing.
\end{alertbox}

\section{The Ten Commands You Will Actually Use}

SQLite's shell has two kinds of input: SQL, which ends with a semicolon, and
\emph{dot-commands}, which do not and which are the shell's own.

\rowcolors{2}{white}{lightbg}
\begin{longtable}{@{}L{4.6cm}L{10.0cm}@{}}
\toprule
\rowcolor{primary!15}
\textbf{Command} & \textbf{What it does} \\
\midrule
\endhead
\texttt{sqlite3 university.db} & Open (or create) a database and enter the
shell \\
\texttt{.tables} & List the tables \\
\texttt{.schema student} & Show the \texttt{CREATE} statement for one table ---
the fastest way to find out what a table actually promises \\
\texttt{.headers on} & Print column names above results. Turn this on first,
every time \\
\texttt{.mode box} & Draw results as a table instead of pipe-separated text \\
\texttt{.read setup.sql} & Run a file of SQL. How every lab in this book starts \\
\texttt{.output results.txt} & Send what follows to a file \\
\texttt{.import data.csv student} & Load a CSV into a table \\
\texttt{.backup backup.db} & Take a consistent copy \emph{while the database is
open} --- \chref{transactions} explains why that is not the same as copying the
file \\
\texttt{.quit} & Leave \\
\bottomrule
\end{longtable}

\begin{examplebox}[The two lines to type at the start of every session]
\begin{lstlisting}[language=bash]
sqlite3 university.db
\end{lstlisting}
then, inside the shell:
\begin{lstlisting}[language=sql]
.headers on
.mode box
PRAGMA foreign_keys = ON;
\end{lstlisting}

\smallskip
The first two make results readable. \textbf{The third is not cosmetic.}
SQLite does not enforce foreign keys unless you ask it to, \emph{every time you
connect}, and \chref{constraints} is largely about what happens when you
forget. Put all three in a file called \texttt{init.sql} and start every
session with \texttt{.read init.sql} until it is a habit.
\end{examplebox}

\section{Resetting to a Known State}

A database is one file, which makes this the simplest reset in the series:

\begin{lstlisting}[language=bash]
# --- 1. Keep a copy before anything risky -------------------------
cp university.db university.db.bak

# --- 2. Start again from nothing ----------------------------------
rm university.db
sqlite3 university.db ".read code/schema/v1_create.sql"

# --- 3. And refill it ----------------------------------------------
python code/schema/seed.py 1000
\end{lstlisting}

\begin{conceptbox}[Do this deliberately once, in Week 2]
Delete your database on purpose, rebuild it from the script, and confirm it is
the same. A student who has done that once is a student who will keep their
schema in a file rather than in the shell's history --- which is the single
habit that decides whether the Week~16 project can be handed to somebody else
and run.
\end{conceptbox}

\section{The Machine Every Timing Was Taken On}

Eleven chapters quote a measured figure. Every one was produced by running the
named file in \texttt{code/} on this machine:

\rowcolors{2}{white}{lightbg}
\begin{longtable}{@{}L{3.6cm}L{11.0cm}@{}}
\toprule
\rowcolor{primary!15}
\textbf{Component} & \textbf{Specification} \\
\midrule
\endhead
Processor & Intel Core i5-6600, 4 cores, 3.30\,GHz \\
Memory & 16\,GB \\
Operating system & Windows 11, 64-bit \\
SQLite & 3.51.2 \\
Python & 3.13.13 \\
Storage & Mechanical disk, not a solid-state drive --- which matters for
\chref{indexes} and is the reason its figures are larger than yours may be \\
\bottomrule
\end{longtable}

\begin{alertbox}[Your numbers will differ, and the ratios should not]
A measurement is a fact about the machine it ran on. If your laptop has a
solid-state disk it will beat these figures, perhaps by a factor of ten on the
chapters that touch storage.

\smallskip
\textbf{What should survive the move is the shape:} the row count at which an
index starts to pay (\chref{indexes}), the direction normalisation moves query
time (\chref{normalforms}), the fact that one writer blocks another
(\chref{transactions}). If a \emph{ratio} in this book does not reproduce on
your machine, that is worth investigating and worth telling your lecturer
about. If an absolute number does not, that is expected.
\end{alertbox}

\section{Scope: What You May Point Any of This At}

\begin{alertbox}[Read this before Chapter 23's laboratory]
Everything in this course runs against \textbf{a database file you created, on
a machine you own}. \chref{project} discusses SQL injection and shows what a
badly built query permits --- against your own university database, with
invented data, in a file you can delete.

\smallskip
\textbf{Running the same thing against any system you do not own and have not
been authorised in writing to test is a criminal offence in Pakistan} under the
\textbf{Prevention of Electronic Crimes Act 2016}, sections~3 and~4, which
cover unauthorised access to an information system and to data.

\smallskip
There is no educational exemption, and a university assignment is not
authorisation. ``I was learning'', ``I did no harm'' and ``the site was
insecure'' are not defences.

\smallskip
In practice this is an easy rule to follow, because the only file named
anywhere in this book is \texttt{university.db} and it is yours.
\end{alertbox}
